\section{Labor backing accounts}
\label{sec:accounts}

This section sets out the paper's method contribution. Labor backing accounts decompose the
whole stock of United States financial claims twice over: once by the income that services each
claim, and once by who ultimately holds, guarantees or owes it. The two decompositions are taken
on the same claim stock in the same period, which is what allows the wage side and the AI side of
the exposure to be compared rather than asserted. Everything later in the paper is built on the
object defined here.

\subsection{Definition and vocabulary}

A financial claim is serviced out of some income stream. A mortgage is serviced out of household
wages; a corporate bond out of business revenue; Treasury debt out of federal receipts, which are
themselves drawn from wages, from capital income and from other sources in measurable
proportions. The \emph{labor backing} of a class of claims is the share of the cash flow that
services it which is labor income. A claim is \emph{directly wage-backed} when wages pay it
without passing through another party's income statement first. It is \emph{indirectly
wage-backed} when wages first become spending, spending becomes business revenue, and that
revenue services the claim.

The distinction is the accounting form of one the rest of the paper uses throughout. The
\emph{first round} is what happens to the party whose own wages fall. The \emph{second round} is
what happens to everyone else, through the spending, house prices and business credit that
follow. Direct backing is the first round seen on the liability side of the economy; indirect
backing is the second.

Two conventions of presentation travel with every result below. \emph{Displacement} is always
stated as a share of the total wage bill, never as a count of jobs and never as a share of one
occupation, because the fiscal and the credit channel both run on dollars of pay. And every
quantitative statement carries a status on its first use. A \emph{replicated} figure was rebuilt
by an instance that had not written the code, from raw data and a published brief alone, inside a
stated tolerance. A \emph{measured} figure comes from the pipeline and clears every applicable
check but has not been independently rebuilt. A \emph{scenario} figure is arithmetic conditional
on a stated assumption. Section~\ref{sec:data} sets out how the labels are assigned and what they
are worth.

The accounts run over thirteen classes of claim drawn from the Federal Reserve financial
accounts, listed with their levels in Table~\ref{tab:classes}. They cover household mortgage and
consumer credit, multifamily and commercial mortgages, Treasury and state and local debt,
corporate bonds and loans, noncorporate business debt, and corporate equity.

\begin{table}[htbp]
\centering
\caption{The thirteen claim classes, levels and labor backing, 2025. Classes marked $\dagger$ are
serviced from business revenue and carry zero direct backing under the one-step rule. Measured;
the home mortgage and multifamily coefficients are replicated.}
\label{tab:classes}
\begin{tabular}{lrrr}
\toprule
Claim class & Level (USD bn) & Labor backing & Wage-backed (USD bn) \\
\midrule
Home mortgage & 13,788.3 & 0.8104 & 11,173.5 \\
Credit card & 1,324.3 & 0.7302 & 967.0 \\
Auto loan & 1,562.2 & 0.7783 & 1,215.9 \\
Student loan & 1,834.7 & 0.8822 & 1,618.5 \\
Other consumer credit & 377.6 & 0.8052 & 304.0 \\
Multifamily mortgage & 2,448.7 & 0.7275 & 1,781.5 \\
Treasury debt & 33,887.1 & 0.6578 & 22,290.5 \\
State and local debt & 3,803.9 & 0.1593 & 606.1 \\
Corporate bonds\,$\dagger$ & 8,075.3 & 0.0000 & 0.0 \\
Corporate loans\,$\dagger$ & 3,903.7 & 0.0000 & 0.0 \\
Noncorporate business debt\,$\dagger$ & 2,453.0 & 0.0000 & 0.0 \\
Commercial mortgage\,$\dagger$ & 3,953.4 & 0.0000 & 0.0 \\
Corporate equity\,$\dagger$ & 71,994.9 & 0.0000 & 0.0 \\
\midrule
All thirteen classes & 149,407.2 & 0.2674 & 39,957.2 \\
\bottomrule
\end{tabular}

\end{table}

\subsection{Two conventions, stated plainly because they are contestable}

\paragraph{Government debt enters by its receipts, not by its obligor.} Treasury debt is not
wage-backed because a household owes it. It is wage-backed to the extent that wage taxes service
it. We therefore assign Treasury debt the labor-linked share of federal receipts, which is social
insurance contributions in full plus the wage share of the individual income tax, measured at
\result{ReceiptsPctLow} percent. State and local debt is treated the same way on its own receipts
mix. The convention is what makes the federal exposure in Section~\ref{sec:holders} an empirical
quantity rather than a definitional one. Without it a government funded entirely from payroll
taxes and a government funded entirely from capital taxes would score identically, which is
plainly the wrong answer for a question about wage risk.

\paragraph{The one-step rule is a convention, and it is the largest judgement in the paper.} We
trace servicing one step and stop. A claim serviced out of business revenue carries zero direct
backing, even though some of that revenue began as somebody's wages. Multifamily mortgages are the
one property class that does not take the zero, because the rent that services them is paid out of
wages directly, with no firm in between.

The rule is the first-round and second-round boundary applied to the liability side, and we adopt
it for that reason rather than because tracing further is impossible. Tracing further is exactly
what the indirect variant does. We report both variants everywhere and never average them. On the
whole claim stock, moving the business classes from zero to their measured indirect share raises
the ratio by \result{OneStepMoveAllClaimsPct} percent, more than every other judgement call in the
accounts put together. A reader who rejects the convention gets a different number, and that
number is printed beside ours rather than buried.

\subsection{The debt-only ratio, and why it is the one we report}

The natural first construction puts every claim in the denominator, corporate equity at market
value included. It has a defect serious enough to disqualify it as a headline. Equity is
\result{EquityShareDenom} percent of the all-claims denominator and moves with asset prices, so
the ratio falls when equity rises and rises when equity falls, independently of anything happening
to wages. It read \result{AllClaims2000} percent at the 2000 equity peak and \result{AllClaims2008}
percent in 2008 after the crash. Read as an exposure indicator it would have called the economy
least exposed in 2000 and most exposed in 2009, which is the wrong sign on both.

The debt-only ratio removes every market-valued equity class from the denominator, leaving claims
carried at par or amortized cost, so the series measures composition rather than valuation. The
exclusion keys on the valuation basis rather than on the name of the class, so a future
equity-like instrument is caught automatically. On the same two dates the debt-only ratio reads
\result{DebtOnly2000} and \result{DebtOnly2008} percent, a three-point move where the all-claims
ratio moved eleven.

On this measure, about \result{DebtOnlyDirect} percent of all US debt, public and private, is
serviced directly from wages, and \result{DebtOnlyIndirect} percent once indirect servicing is
counted. The corresponding all-claims figures are \result{AllClaimsDirect} and
\result{AllClaimsIndirect} percent. All four are replicated: an independent rebuild reproduced
each of them to within 0.001 (Section~\ref{sec:replication}).

Two further properties favour the debt-only construction, and both are replicated. Over 1952 to
2025 its coefficient of variation is \result{DebtOnlyCV} against \result{AllClaimsCV} for the
all-claims series, so it is roughly half as volatile. And it is far less sensitive to the one-step
rule, which moves it by \result{OneStepMovePct} percent against
\result{OneStepMoveAllClaimsPct} percent for the all-claims ratio. The statistic we put forward is
the one that is stable and the one that depends least on our own largest convention. That is not a
coincidence, it is the reason for the construction.

\subsection{Sensitivity}
\label{sec:sensitivity}

Table~\ref{tab:sensitivity} moves each named judgement call one at a time and reports the
resulting ratio. The one-step rule leads it, as it leads everything in this section, and nothing
else moves the debt-only ratio by more than ten percent. The table is one at a time and so
understates the true range if the calls interact, which the one-step rule does with all of them.

\begin{table}[htbp]
\centering
\caption{Sensitivity of the debt-only labor backing ratio, one judgement call at a time.
Measured.}
\label{tab:sensitivity}
\begin{tabular}{lrr}
\toprule
Judgement call & Debt-only ratio & Move (per cent) \\
\midrule
Central case & 0.5216 & +0.00 \\
The one-step rule: business classes take the indirect share & 0.6020 & +15.40 \\
federal receipts: the 77.7 percent upper bound & 0.5737 & +9.99 \\
commercial mortgage treated like multifamily & 0.5588 & +7.12 \\
mortgage backing: the superseded A38 definition & 0.5150 & -1.27 \\
state and local: all personal current taxes & 0.5255 & +0.75 \\
mortgage backing: SIPP balances & 0.5181 & -0.67 \\
rent backing: the A38 definition & 0.5207 & -0.17 \\
other consumer: the card share alone & 0.5213 & -0.07 \\
\bottomrule
\end{tabular}

\end{table}

One bound travels with every fiscal magnitude in the paper. The labor-linked share of federal
receipts is built from the wage share of adjusted gross income, and realized capital gains sit
inside that base while bearing preferential rates. On Internal Revenue Service Statistics of
Income data for 2021 to 2023, realized gains run from \result{GainsAGILow} to
\result{GainsAGIHigh} percent of adjusted gross income, so the labor-linked receipts share is a
band, \result{ReceiptsPctLow} to \result{ReceiptsPctHigh} percent, rather than a point. We publish
the lower end throughout. The federal exposure share moves by less than one percent across the
band. The fiscal magnitudes in Section~\ref{sec:displacement} are more sensitive to it, and we
have not bounded that.

\subsection{Who is exposed, and how}
\label{sec:holders}

The second decomposition asks who holds each claim. Table~\ref{tab:holders} gives the holder
shares of the directly wage-backed stock. The federal government is the largest single holder or
guarantor, at \result{SovereignCreditor} percent, through agency mortgage guarantees, the student
loan book and central bank holdings. Banks and the rest of the world come next, at roughly
eighteen percent each.

\begin{table}[htbp]
\centering
\caption{Holders of directly wage-backed claims, 2025. Measured.}
\label{tab:holders}
\begin{tabular}{lrr}
\toprule
Holder & Wage-backed claims held (USD bn) & Share \\
\midrule
Federal government & 12,712.5 & 0.3182 \\
Banks & 7,329.8 & 0.1834 \\
Rest of the world & 7,271.1 & 0.1820 \\
Other financial & 6,158.8 & 0.1541 \\
Households, direct & 2,449.6 & 0.0613 \\
State and local government & 1,863.5 & 0.0466 \\
Insurers & 816.2 & 0.0204 \\
Pension funds & 444.5 & 0.0111 \\
Nonfinancial business & 249.7 & 0.0062 \\
Unallocated remainder & 661.4 & 0.0166 \\
\bottomrule
\end{tabular}

\end{table}

Holding is not the only way a claim reaches a balance sheet. For an agency-guaranteed mortgage the
federal government bears a credit loss. For Treasury debt serviced from wage taxes it bears a
revenue shortfall on a claim it owes. Both are ways a fall in wages reaches the same institution,
and a measure that counted only the first would understate the concentration by a factor of more
than two. We therefore report the federal position as a union of three roles, holder, guarantor
and debtor, and state it that way wherever it appears.

\begin{quote}
The federal government is exposed, as holder, guarantor or debtor, on about
\result{SovereignUnion} percent of directly wage-backed claims: \result{SovereignCreditor} percent
as creditor or guarantor and \result{SovereignDebtor} percent as the debtor on debt serviced from
wage taxes, less the \result{OverlapBn} billion dollar overlap where the federal sector holds its
own debt. Measured, with an agreed range of \result{AgreedRangeLow} to \result{AgreedRangeHigh}
percent after independent rebuild.
\end{quote}

The union is not a portfolio share and must not be read as one. A reader who takes
\result{SovereignUnion} percent to mean that the government owns four fifths of the claim stock
has it wrong. It means that four fifths of the claims wages pay reach the federal balance sheet
through one of three doors.

The figure is stable against every measurement call and unstable against two structural ones, and
the honest way to report it is to say both in the same breath. Varying the Treasury series
definition, the mortgage backing definition, the rent backing definition, the state and local
split, the consumer credit blend, or whether the central bank counts as federal moves the share by
at most \result{MeasurementMaxMovePct} percent, and the central bank question moves it by exactly
zero. The structural choices in Table~\ref{tab:structural} move it a great deal more. Counting
indirectly wage-backed claims takes it to \result{StructIndirect} percent; not treating agency
pools as federal takes it to \result{StructNoPools} percent; excluding the debtor position
entirely leaves \result{StructObligorExcluded} percent. This is a stable measurement under a
stated convention, and the convention does heavy lifting.

\begin{table}[htbp]
\centering
\caption{Federal exposure under alternative structural choices. Measured.}
\label{tab:structural}
\begin{tabular}{lrr}
\toprule
Structural choice & Federal exposure & Move (per cent) \\
\midrule
Central case & 0.7936 & 0.00 \\
Commercial mortgage treated as rent-serviced & 0.7408 & -6.65 \\
Indirectly wage-backed claims included & 0.4517 & -43.08 \\
Agency pools not treated as federal & 0.5865 & -26.09 \\
Federal debtor position excluded & 0.3215 & -59.49 \\
\bottomrule
\end{tabular}

\end{table}

\subsection{The series, 1952 to 2025}

Table~\ref{tab:series} runs the accounts back to 1952. Two things in it move differently, and
reporting the series from 1970 would misdescribe both.

The ratio itself is close to flat across seventy-three years. The all-claims measure stays between
\result{AllClaimsSeriesLow} and \result{AllClaimsSeriesHigh} percent with no trend, and the
debt-only measure between \result{DebtOnlySeriesLow} and \result{DebtOnlySeriesHigh} percent. The
share of the national debt stock that wages service is not a recent development.

The federal share of it traces a U. It stands at \result{SovUnion1952} percent in 1952, falls to
\result{SovUnion1970} percent in 1970, and climbs back to \result{SovUnion2025} percent in 2025.
The 1952 level is almost entirely wartime Treasury debt: the debtor position alone is
\result{SovObligor1952} percent that year against \result{SovHeld1952} percent held or guaranteed.
It falls to 1970 as that debt shrinks against a growing claim stock. It then climbs in two datable
steps. The agency book grows between 1970 and 1990, taking the held-or-guaranteed position from
\result{SovHeld1970} to \result{SovHeld1990} percent. Federal debt then grows between 2008 and
2020, taking the debtor position from \result{SovObligor2008} to \result{SovObligor2020} percent.

Said plainly, because it qualifies the headline: a large part of the recent rise is growth in
federal debt itself rather than new exposure to households. The held-or-guaranteed position peaked
at \result{SovHeld2020} percent in 2020 and has since fallen to \result{SovHeld2025} percent,
while the whole net increase since 2008 sits in the debtor position. A reader who wants to know
whether the government has been writing more wage-backed credit should look at the creditor
column, where the answer for the last five years is no.

\begin{table}[htbp]
\centering
\caption{Labor backing and federal exposure, selected years, 1952 to 2025. The creditor and debtor
columns overlap where the federal sector holds its own debt, so they do not sum to the union.
Measured.}
\label{tab:series}
\begin{tabular}{lrrrrr}
\toprule
Year & Debt-only ratio & All-claims ratio & Federal, union & \quad as creditor & \quad as debtor \\
\midrule
1952 & 0.5208 & 0.2837 & 0.5227 & 0.0853 & 0.5102 \\
1970 & 0.4520 & 0.2767 & 0.3356 & 0.1204 & 0.2965 \\
1990 & 0.4600 & 0.3549 & 0.5638 & 0.2482 & 0.3647 \\
2000 & 0.4687 & 0.2673 & 0.5617 & 0.3154 & 0.3043 \\
2008 & 0.5006 & 0.3763 & 0.5582 & 0.2977 & 0.2859 \\
2020 & 0.5037 & 0.2907 & 0.7831 & 0.3923 & 0.5225 \\
2025 & 0.5216 & 0.2703 & 0.7936 & 0.3215 & 0.5520 \\
\bottomrule
\end{tabular}

\end{table}

\subsection{The other side of the bet}

The same decomposition applied to claims on AI capital gives the comparison the paper is for. The
federal government holds about \result{AIFederalShareRounded} percent of them, across a threefold
variation in the assumed scale of AI capital, \result{AIFederalShareLow} to
\result{AIFederalShareHigh} percent, with a central assumption that AI capital is
\result{AIEquityScale} percent of US nonfinancial corporate equity. That scale is a scenario
parameter and is never asserted: no filing discloses an AI-attributable market value, so we sweep
it and report the whole sweep.

One caveat is load-bearing and belongs in the same sentence as the number. The AI side is built
from nine SEC registrants' on-balance-sheet facts. Off-balance-sheet vehicles, which the Bank for
International Settlements identifies as dominant in data centre financing, contribute nothing to
it, and neither do GPU-backed lending or data centre securitizations. Every level on the AI side
is therefore a lower bound, and the federal share of it is correspondingly an upper bound: adding
the unmeasured financing would enlarge the denominator and push the share below
\result{AIFederalShareRounded} percent. The gap between the two sides, about \result{HolderGap}
percentage points, is measured on the wage side and bounded on the AI side, and it would widen
rather than narrow if the missing financing were found.

\subsection{Households, by position in the wage distribution}

The accounts also split the household part of the wage-backed stock by wage quintile, in
Table~\ref{tab:quintile}. The gradient is the finding. Wage-backed claims per dollar of wages fall
as pay rises, from \result{QTwoPerDollar} in the second quintile to \result{QTopPerDollar} at the
top. The top quintile earns \result{QTopWageShare} percent of the wage bill and carries
\result{QTopClaimShare} percent of the household claims, so the claims are spread far more evenly
than the income that services them. That is the sense in which a wage shock is regressive before
any policy response has been chosen: the same proportional fall in pay consumes a larger share of
the servicing capacity of households further down.

The magnitude for the bottom quintile is withdrawn and does not appear in
Table~\ref{tab:quintile}. An independent rebuild landed a factor of 2.5 to 3.4 away from ours, and
we cannot show its reading was wrong, because our own brief defined neither the universe nor the
normalisation of that cell. The direction of the gradient stands and was reproduced; the
bottom-quintile level is reported nowhere in this paper. Section~\ref{sec:limitations} lists it
among the withdrawn quantities.

\begin{table}[htbp]
\centering
\caption{Wage-backed household claims by wage quintile. The bottom-quintile magnitude is
withdrawn. Measured.}
\label{tab:quintile}
\begin{tabular}{lrrrr}
\toprule
Wage quintile & Share of the wage bill & Share of household wage-backed claims & Claims per wage dollar & Same, own-survey denominator \\
\midrule
Bottom quintile & 0.0324 & 0.1344 & 4.0359\,$\dagger$ & 4.0257\,$\dagger$ \\
Second quintile & 0.0895 & 0.1263 & 1.3728 & 1.4364 \\
Middle quintile & 0.1426 & 0.1716 & 1.1703 & 1.2533 \\
Fourth quintile & 0.2233 & 0.2323 & 1.0120 & 1.1088 \\
Top quintile & 0.5122 & 0.3354 & 0.6370 & 0.5986 \\
\bottomrule
\end{tabular}

\end{table}

\subsection{What has been independently rebuilt}
\label{sec:replication}

Three rebuild rounds have been run against these accounts, and what each established is worth
stating precisely, because the available words carry different weights.

Round two scored 127 quantities, 79 of them against sealed counterparts, by an instance that read
no project code. Within this section it reproduced the home mortgage and multifamily backing
coefficients to the last published digit, rebuilt from survey mortgage service and gross rent, and
reached a federal exposure share of 78 percent against our \result{SovereignUnion} percent, a
difference that decomposes almost entirely to the definition of the Treasury class. Round three
took the debt-only ratio and the capital tax assembly of Section~\ref{sec:fiscal}, the two
constructions the earlier rounds did not cover, and matched all 32 sealed quantities inside
tolerance, 30 of them within rounding. That is the basis on which the debt-only pair, the equity
share of the denominator, both coefficients of variation and the size of the one-step move are
labelled replicated here.

The protocol carries one qualification, stated here and again in Section~\ref{sec:data}. The blind
was procedural rather than enforced: the rebuilder worked on the same machine, with the sealed
file reachable at a path the brief itself names, and reported it unopened until the rebuild was
final. The claim is therefore that these quantities were independently rebuilt under a reported
blind protocol, not that they were independently verified.
